% Generated draft: requires mathematical and source review.
\documentclass[11pt]{article}
\usepackage[margin=25mm]{geometry}
\usepackage{fontspec}
\setmainfont{Latin Modern Roman}
\usepackage{amsmath,amssymb,mathtools,unicode-math}
\setmathfont{Latin Modern Math}
\usepackage{enumitem,needspace,xurl,hyperref}
\setlength{\parindent}{0pt}
\setlength{\parskip}{5pt}
\setlength{\emergencystretch}{3em}
\newcommand{\N}{\mathbb{N}}
\newcommand{\Z}{\mathbb{Z}}
\newcommand{\Q}{\mathbb{Q}}
\newcommand{\R}{\mathbb{R}}
\newcommand{\C}{\mathbb{C}}
\newcommand{\F}{\mathbb{F}}
\newcommand{\A}{\mathbb{A}}
\newcommand{\PP}{\mathbb{P}}
\newcommand{\E}{\mathbb{E}}
\newcommand{\eps}{\varepsilon}
\newcommand{\dd}{\,\mathrm{d}}
\begin{document}
\section*{UnsolvedMath ID: 2796}
\textbf{Source problem code:} catalogue entry
\subsection{Definitions and mathematical statement}
Let $S$ be a closed connected oriented surface of genus $g\ge1$. Equip it with a complex structure, so
\[
 M=\operatorname{Sym}^g(S)=S^g/\mathfrak S_g
\]
is the smooth complex $g$-fold parametrizing unordered $g$-tuples with multiplicity. Write $\operatorname{Mod}(S)=\pi_0\operatorname{Diff}^+(S)$ and $\pi_0\operatorname{Diff}^+(M)$ for orientation-preserving smooth mapping class groups. The source invokes Perutz's construction of symplectic mapping-class actions for suitable forms, giving the induced smooth-class homomorphism
\[
 \rho_g:\operatorname{Mod}(S)\longrightarrow\pi_0\operatorname{Diff}^+(M).
\]
This is not defined by simply declaring the obvious coordinatewise quotient map smooth. The source explicitly warns that the map $\{x_i\}\mapsto\{\varphi(x_i)\}$ for a surface diffeomorphism is generally only a homeomorphism in the complex-induced smooth structure; it is holomorphic when $\varphi$ is holomorphic for the chosen structure.

Part (a) asks for
\[
 \ker\rho_g=\{[\varphi]:\text{the induced diffeomorphism class of $M$ is smoothly isotopic to the identity}\},
\]
and for its based analogue. For a marked point $z\in S$ the natural source group is $\operatorname{Mod}(S,z)$, and the background's symplectic comparison uses $M_z=\operatorname{Sym}^g(S\setminus\{z\})$. A corresponding smooth kernel would depend on the exact based target and allowed isotopies at the noncompact ends. The statement's phrase ``based case'' does not specify these conventions or whether a marked closed target is meant, so no single unqualified based smooth target is substituted.

For (b), let $\alpha_1,\ldots,\alpha_g$ be pairwise disjoint embedded circles in $S$, and similarly let $\beta_1,\ldots,\beta_g$ be such circles. The quotient map restricts to an embedding of their product, giving unparametrized embedded tori
\[
 T_\alpha=\operatorname{image}(\alpha_1\times\cdots\times\alpha_g\to M),\qquad
 T_\beta=\operatorname{image}(\beta_1\times\cdots\times\beta_g\to M).
\]
The question asks for necessary and sufficient conditions on the curve systems for smooth isotopy of these tori. Smooth isotopy means a family of embeddings of a $g$-torus with those endpoint images, with no prescribed endpoint parametrization. It separately asks when there is a Lagrangian isotopy for an appropriate fixed symplectic form $\omega$ on $M$: all intermediate tori must have dimension $g$ and zero pullback of $\omega$. Hamiltonian isotopy is a stronger condition, generated by a time-dependent Hamiltonian flow. The source does not specify the full permitted class of forms for this existential choice; it cannot be replaced by a claim for every form. The displayed curve hypothesis is disjointness; no unstated cut-system or homological-independence condition is added to the general question.

The supplied partial results and qualifications are:
\begin{itemize}
\item Perutz constructs the relevant symplectic actions. Clarkson proves injectivity of the based mapping class group into symplectomorphisms of $M_z$ modulo Hamiltonian diffeomorphisms. This faithfulness is not faithfulness in the smooth mapping class group: being smoothly isotopic to the identity need not mean Hamiltonian. Thus it does not determine the closed or based smooth kernel.
\item For appropriate punctured symplectic forms, the relevant Heegaard tori are exact, as used by Hendricks and Hendricks--Lidman--Lipshitz, and Heegaard Floer homology agrees with ordinary Lagrangian intersection Floer homology in that setting. These exactness/form hypotheses are not asserted for arbitrary closed-product tori.
\item A handleslide replacing one attaching curve $\alpha_1$ by $\alpha'_1$, with the other curves retained, produces smoothly isotopic product tori and, for suitable forms, Hamiltonian-isotopic tori by Perutz's handleslide theorem. This settles that move, not an if-and-only-if classification of arbitrary systems.
\item For $g>1$, $\pi_1(M)\cong H_1(S;\Z)$. The Torelli group, the kernel of the action on $H_1(S;\Z)$, acts trivially on this fundamental group and on the homology of $M$. This does not imply trivial action in the smooth mapping class group.
\item In genus $2$, the Abel--Jacobi map presents $\operatorname{Sym}^2(S)$ as a blow-up of the Jacobian torus $H^1(S;\R)/H^1(S;\Z)$. The source suggests that trivial Torelli action on the Jacobian might help show smooth triviality, and speculates similarly about Torelli or the Johnson kernel in higher genus. These are suggested approaches, not established kernel descriptions. The source's occasional $\Sigma$ denotes the same surface in this discussion.
\end{itemize}
The source relates the question to whether Heegaard Floer theory reflects symplectic information beyond smooth isotopy of the tori and cites Seidel for related smooth-versus-symplectic phenomena. The literature triage checked 17 August 2026 labels partial progress: the based symplectic and handleslide cases are known, but the smooth kernels and general smooth/Lagrangian torus-isotopy classifications remain open. The main registry label is open. The genus-zero empty symmetric product is a separate degenerate case rather than the setting of these positive-genus Heegaard tori.
\subsection{Short English statement}
Determine which surface mapping classes become smoothly trivial on a symmetric product, and when product tori from curve systems are smoothly or Lagrangian isotopic. Known based symplectic faithfulness and Hamiltonian handleslides do not determine the smooth kernels or the full isotopy classification.
\subsection{Sources}
\begin{itemize}
\item Local UnsolvedMath registry, permanent ID $2796$, exact source code \texttt{catalogue entry}, ``Kirby Problem 2.48,'' including its literature triage checked 17 August 2026.
\item Baykur, Kirby and Ruberman, \emph{K3---A New Problem List in Low-Dimensional Topology} (2026), Problem 2.48, \url{https://math.berkeley.edu/sites/default/files/surv-295-ruberman-watermarked-author-pdf.pdf}.
\item [Sei97] Paul Seidel, \emph{Floer homology and the symplectic isotopy problem}, PhD thesis, University of Oxford, 1997.
\item [Per08a] Tim Perutz, \emph{Lagrangian matching invariants for fibred four-manifolds. II}, Geometry \& Topology \textbf{12}(3) (2008), 1461--1542, DOI 10.2140/gt.2008.12.1461.
\item [Cla17] Corrin Clarkson, \emph{Three-manifold mutations detected by Heegaard Floer homology}, Algebraic \& Geometric Topology \textbf{17}(1) (2017), 1--16, \url{https://doi.org/10.2140/agt.2017.17.1}.
\item [Per08b] Timothy Perutz, \emph{Hamiltonian handleslides for Heegaard Floer homology}, Proceedings of G\"okova Geometry-Topology Conference 2007, 15--35, G\"okova Geometry/Topology Conference (2008), \url{https://gokovagt.org/proceedings/2007/2.pdf}.
\item [Hen12] Kristen Hendricks, \emph{A rank inequality for the knot Floer homology of double branched covers}, Algebraic \& Geometric Topology \textbf{12}(4) (2012), 2127--2178, DOI 10.2140/agt.2012.12.2127.
\item [HLL22a] Kristen Hendricks, Tye Lidman and Robert Lipshitz, \emph{Rank inequalities for the Heegaard Floer homology of branched covers}, Documenta Mathematica \textbf{27} (2022), 581--612.
\end{itemize}
\end{document}
